\section{Primary descent over multiplicity strata}
\label{sec:primary-descent}
The following descent statement concerns ideals, not just the
permutation action on their generic points.

\begin{lemma}[Etale descent of weighted prime divisors]
\label{lem:etale-primary-descent}
Let $X$ be a regular integral Noetherian $\C$-scheme and
$p:\widetilde X\to X$ a finite etale Galois cover, with group $G$.
Let $\{D_\alpha\}$ be a finite $G$-stable collection of distinct
integral effective Cartier divisors on $\widetilde X$. Work
componentwise if the cover is disconnected. Suppose the positive
integers $b_\alpha$ are constant on $G$-orbits. For an orbit $O$,
there is a unique integral effective Cartier divisor $D_O$ on $X$
whose ideal satisfies
\[
 p^*\mathcal I_{D_O}=\prod_{\alpha\in O}\mathcal I_{D_\alpha}
                      =\bigcap_{\alpha\in O}\mathcal I_{D_\alpha}.
\]
If a descended ideal $\mathcal J$ pulls back to
$\prod_\alpha\mathcal I_{D_\alpha}^{b_\alpha}$, then its complete
primary decomposition is
\begin{equation}\label{eq:etale-full-primary}
 \mathcal J=\prod_O\mathcal I_{D_O}^{b_O}
            =\bigcap_O\mathcal I_{D_O}^{b_O}.
\end{equation}
There are no embedded associated points, and the formula holds at
intersections as well as generically. The descended prime divisor
need not be smooth or geometrically unibranched.
\end{lemma}
\begin{proof}
In each regular local ring upstairs, equations for distinct prime
Cartier divisors are distinct prime factors. Unique factorization
gives product equal to intersection, with or without positive
powers. The reduced orbit union has an invariant invertible ideal
subsheaf of $\mathcal O_{\widetilde X}$ with its canonical descent
data. Faithfully flat descent gives an invertible ideal subsheaf
of $\mathcal O_X$. It defines an effective Cartier divisor. Its
support is the image of any one member of the orbit, hence is
irreducible, since the map is finite. Reducedness descends
faithfully flat, so this divisor is integral.

The same factorization upstairs, and flatness for finite
intersections, identify the pullbacks of all three ideals in
\eqref{eq:etale-full-primary}. Faithfulness makes them equal
downstairs. Locally downstairs each $\mathcal I_{D_O}$ is generated
by a prime element of a regular local ring; its $b_O$th power is
primary. Their intersection is irredundant after discarding
components absent from the chart, so it is the complete primary
decomposition and has no embedded primes. No normal-crossing
hypothesis was used. An orbit can have several members through a
point upstairs, giving several etale branches of the integral
descended divisor. This explains the final distinction.
\end{proof}

Apply the lemma to the relative generating pencil Grassmannian over
a multiplicity stratum. Ordering the distinct support points gives
a finite etale cover with permutations among equal multiplicities.
On it the divisors $T_i,E_{ij}$ are the relative rank-one Schubert
divisors. Their ideals are prime Cartier, and their weights are
$\binom{d_i}2,d_id_j$. Formula~\eqref{eq:etale-full-primary} supplies
the whole descended primary decomposition. In particular it remains
valid at a braid intersection where several values agree. The
fixed labelled divisors are smooth; the descended orbit divisors
are asserted integral, not automatically smooth or unibranched.

This argument applies to the earlier pencil calculation. For the
contact-hyperplane family, Theorem~\ref{thm:universal-hyperplane}
gives a scheme isomorphism with the quotient incidence and does not
require a conjectural primary descent rule.
